% (c) Nikita Lisitsa, lisyarus@gmail.com, 2026

\documentclass[10pt]{beamer}

\usepackage[T2A]{fontenc}
\usepackage[russian]{babel}
\usepackage{minted}

\usepackage[most]{tcolorbox}
\tcbuselibrary{minted}

\usepackage{graphicx}
\graphicspath{ {./images/} }

\usepackage{adjustbox}

\usepackage{color}
\usepackage{soul}

\usepackage{hyperref}

\usetheme{metropolis}

\definecolor{red}{rgb}{1,0,0}
\definecolor{codebg}{RGB}{29,35,49}
\setminted{fontsize=\footnotesize}

\makeatletter
\newcommand{\slideimage}[1]{
  \begin{figure}
    \begin{adjustbox}{width=\textwidth, totalheight=\textheight-2\baselineskip-2\baselineskip,keepaspectratio}
      \includegraphics{#1}
    \end{adjustbox}
  \end{figure}
}
\makeatother

\title{Язык программирования C++}
\subtitle{Лекция 12: наследование, полиморфизм, виртуальные функции}
\date{2026}
\author{Никита Лисица}

\setbeamertemplate{footline}[frame number]

\begin{document}

\frame{\titlepage}

\begin{frame}[fragile]
\frametitle{Напоминание про наследование}
\begin{itemize}
\item Позволяет нам реализовать один класс поверх другого, переиспользуя его поля и методы:
\usemintedstyle{lightbulb}
\begin{tcblisting}{listing engine=minted, minted language=cpp, minted options={fontsize=\scriptsize}, listing only, colback=codebg, colframe=codebg, rounded corners}
class Button {
protected:
  const char* caption_;
  Action clickAction_;

public:
  ...

  virtual void draw() {
    drawRectangle(Color::Gray);
    drawText(caption_);
  }

  virtual void click() {
    clickAction_.execute();
  }
};
\end{tcblisting}
\end{itemize}
\end{frame}

\begin{frame}[fragile]
\frametitle{Напоминание про наследование}
\begin{itemize}
\item Позволяет нам реализовать один класс поверх другого, переиспользуя его поля и методы:
\usemintedstyle{lightbulb}
\begin{tcblisting}{listing engine=minted, minted language=cpp, minted options={fontsize=\scriptsize}, listing only, colback=codebg, colframe=codebg, rounded corners}
class ColoredButton : public Button {
protected:
  Color color_;

public:
  ...

  // Автоматически virtual
  void draw() override {
    drawRectangle(color_);
    drawText(caption_);
  }

  // Автоматически virtual
  void click() override {
    color_ = Color::Red;
    Button::click();
  }
};
\end{tcblisting}
\end{itemize}
\end{frame}

\begin{frame}[fragile]
\frametitle{Override}
\begin{itemize}
\item Допустим, мы перепутали название метода в классе-наследнике:
\usemintedstyle{lightbulb}
\begin{tcblisting}{listing engine=minted, minted language=cpp, minted options={fontsize=\scriptsize}, listing only, colback=codebg, colframe=codebg, rounded corners}
class ColoredButton : public Button {
public:
  void drawButton() { ... }
  ...
};

Button * button = new ColoredButton(...);
button->draw(); // Что произойдёт?
\end{tcblisting}
\pause
\usemintedstyle{tango}
\item Вызовется метод \mintinline{cpp}|Button::draw()|, так как метод \mintinline{cpp}|ColoredButton::drawButton()| не имеет к нему никакого отношения
\end{itemize}
\end{frame}

\begin{frame}[fragile]
\frametitle{Override}
\begin{itemize}
\usemintedstyle{tango}
\item Ключевое слово \mintinline{cpp}|override| заставит компилятор проверить, что метод \textit{переопределяет} какой-то вирутальный метод базового класса
\pause
\usemintedstyle{lightbulb}
\begin{tcblisting}{listing engine=minted, minted language=cpp, minted options={fontsize=\scriptsize}, listing only, colback=codebg, colframe=codebg, rounded corners}
class ColoredButton : public Button {
public:
  // Ошибка компиляции
  void drawButton() override { ... }
  ...
};
\end{tcblisting}
\pause
\usemintedstyle{tango}
\item Хорошая идея -- \textbf{всегда} писать \mintinline{cpp}|override|, когда вы переопределяете виртуальный метод
\end{itemize}
\end{frame}

\begin{frame}[fragile]
\frametitle{Override vs overload}
\begin{itemize}
\usemintedstyle{tango}
\item Переопределение методов (override) не стоит путать с перегрузкой (overload) функций или методов
\pause
\item Переопределение: метод в классе-наследнике перекрывает виртуальный метод базового класса:
\usemintedstyle{lightbulb}
\begin{tcblisting}{listing engine=minted, minted language=cpp, minted options={fontsize=\scriptsize}, listing only, colback=codebg, colframe=codebg, rounded corners}
class B : public A {
public:
  void exec() override { ... }
};
\end{tcblisting}
\pause
\usemintedstyle{tango}
\item Перегрузка: две функции (или два метода одного класса) имеют одинаковое имя, но разные аргументы:
\usemintedstyle{lightbulb}
\begin{tcblisting}{listing engine=minted, minted language=cpp, minted options={fontsize=\scriptsize}, listing only, colback=codebg, colframe=codebg, rounded corners}
class B {
public:
  void exec(int x) { ... }
  void exec(float x) { ... }
};
\end{tcblisting}
\end{itemize}
\end{frame}

\begin{frame}[fragile]
\frametitle{Пример: странный связный список, list.h}
\begin{itemize}
\usemintedstyle{lightbulb}
\begin{tcblisting}{listing engine=minted, minted language=cpp, minted options={fontsize=\scriptsize}, listing only, colback=codebg, colframe=codebg, rounded corners}
class List {
private:
  int value;
  List* next;

public:
  List(int value);
  ~List();

  void push_back(int value);
  size_t length() const;
};
\end{tcblisting}
\end{itemize}
\end{frame}

\begin{frame}[fragile]
\frametitle{Пример: странный связный список, list.cpp}
\begin{itemize}
\usemintedstyle{lightbulb}
\begin{tcblisting}{listing engine=minted, minted language=cpp, minted options={fontsize=\scriptsize}, listing only, colback=codebg, colframe=codebg, rounded corners}
#include "list.h"

List::List(int value) {
  this->value = value; this->next = nullptr;
}

~List::List() { /* Упражнение :) */ }

void List::push_back(int value) {
  List* current = this;
  while (current->next != NULL)
    current = current->next;
  current->next = new List(value);
}

size_t List::length() const {
  // В константном методе this имеет тип const List* const
  size_t count = 0; const List* current = this;
  while (current->next != NULL) {
    current = current->next; ++count;
  }
  return count;
}
\end{tcblisting}
\end{itemize}
\end{frame}

\begin{frame}[fragile]
\frametitle{Пример: странный связный список}
\begin{itemize}
\usemintedstyle{tango}
\item Почему он странный?
\pause
\item Смешаны понятия списка и одного узла (node)
\pause
\item Не бывает пустым
\pause
\item Все операции за линейное время
\pause
\item Это не страшно, пример педагогический :)
\end{itemize}
\end{frame}

\begin{frame}[fragile]
\frametitle{Пример: странный двусвязный список, double\_list.h}
\begin{itemize}
\usemintedstyle{tango}
\item Позже решили добавить двусвязный список, отнаследовшись от уже имеющегося -- многие операции вообще не изменятся: \mintinline{cpp}|push_back(), length()|
\pause
\usemintedstyle{lightbulb}
\begin{tcblisting}{listing engine=minted, minted language=cpp, minted options={fontsize=\scriptsize}, listing only, colback=codebg, colframe=codebg, rounded corners}
#include "list.h"

class DoubleList : public List {
private:
  DoubleList* prev;

public:
  DoubleList(int val);
  ~DoubleList();
  void push_back(int val);
  void pop_back(); // Новый метод!
};
\end{tcblisting}
\end{itemize}
\end{frame}

\begin{frame}[fragile]
\frametitle{Пример: странный двусвязный список, double\_list.cpp}
\begin{itemize}
\usemintedstyle{lightbulb}
\begin{tcblisting}{listing engine=minted, minted language=cpp, minted options={fontsize=\scriptsize}, listing only, colback=codebg, colframe=codebg, rounded corners}
#include "double_list.h"

DoubleList::DoubleList(int value)
  : List(value), prev(nullptr)
{}

~DoubleList::DoubleList() { /* Упражнение :) */ }

void DoubleList::push_back(int value) {
  DoubleList* current = this;
  while (current->next != NULL)
    current = current->next;
  current->next = new DoubleList(value);
  current->next->prev = current;
}

void DoubleList::pop_back() { /* Упражнение :) */ }
\end{tcblisting}
\end{itemize}
\end{frame}

\begin{frame}[fragile]
\frametitle{Пример: странный двусвязный список}
\begin{itemize}
\usemintedstyle{tango}
\item В примере есть три ошибки:
\pause
\begin{enumerate}
\item Поля \mintinline{cpp}|class List| должны быть не \mintinline{cpp}|private|, а \mintinline{cpp}|protected|
\pause
\item У метода \mintinline{cpp}|push_back()| не хватает \mintinline{cpp}|virtual| и \mintinline{cpp}|override|
\end{enumerate}
\pause
\usemintedstyle{lightbulb}
\begin{tcblisting}{listing engine=minted, minted language=cpp, minted options={fontsize=\scriptsize}, listing only, colback=codebg, colframe=codebg, rounded corners}
class List {
protected:
  int val;
  List *next;

public:
  virtual void push_back(int val);
};

class DoubleList : public List {
...
public:
  void push_back(int val) override;
};
\end{tcblisting}
\end{itemize}
\end{frame}

\begin{frame}[fragile]
\frametitle{Пример: странный двусвязный список}
\begin{itemize}
\usemintedstyle{tango}
\item В примере есть три ошибки:
\pause
\begin{enumerate}
\setcounter{enumi}{2}
\item В методе \mintinline{cpp}|DoubleList::push_back()| некорректное приведение типов \mintinline{cpp}|List* -> DoubleList*|
\end{enumerate}
\pause
\usemintedstyle{lightbulb}
\begin{tcblisting}{listing engine=minted, minted language=cpp, minted options={fontsize=\scriptsize}, listing only, colback=codebg, colframe=codebg, rounded corners}
// Ошибка: приведение List* к DoubleList*
current = current->next;

// ОК: приведение DoubleList* к List*
сurrent->next = new DoubleList(value);

// Ошибка: сurrent->next имеет тип List*, у него
// нет поля prev
сurrent->next->prev = current;
\end{tcblisting}
\end{itemize}
\end{frame}

\begin{frame}[fragile]
\frametitle{Приведение типов при наследовании}
\begin{itemize}
\usemintedstyle{tango}
\item При публичном наследовании класс-наследник `умеет' всё, что умеет базовый класс, содержит тот же публичный интерфейс (поля и методы) \pause \(\Longrightarrow\) указатель на класс-наследник неявно приводится к указателю на базовый класс
\pause
\item Обратное, вообще говоря, неверно (например, у \mintinline{cpp}|List| нет поля \mintinline{cpp}|prev| и метода \mintinline{cpp}|pop_back|)
\pause
\item Если мы точно знаем, что по данному указателю лежит конкретный класс-наследник, можно сделать \textit{явное приведение типов}:
\usemintedstyle{lightbulb}
\begin{tcblisting}{listing engine=minted, minted language=cpp, minted options={fontsize=\scriptsize}, listing only, colback=codebg, colframe=codebg, rounded corners}
List* list = ...;

// Знаем, что там на самом деле DoubleList:
DoubleList* dlist = (DoubleList*)list;
dlist->pop_back();
\end{tcblisting}
\end{itemize}
\end{frame}

\begin{frame}[fragile]
\frametitle{Исправление ошибки}
\begin{itemize}
\usemintedstyle{lightbulb}
\begin{tcblisting}{listing engine=minted, minted language=cpp, minted options={fontsize=\scriptsize}, listing only, colback=codebg, colframe=codebg, rounded corners}
void DoubleList::push_back(int value) {
  DoubleList* current = this;
  while (current->next != NULL)
    current = (DoubleList*)current->next;
  current->next = new DoubleList(value);
  ((DoubleList*)current->next)->prev = current;
}
\end{tcblisting}
\end{itemize}
\end{frame}

\begin{frame}[fragile]
\frametitle{Наследование: плюсы}
\begin{itemize}
\item Получается гибкий, расширяемый код
\usemintedstyle{lightbulb}
\begin{tcblisting}{listing engine=minted, minted language=cpp, minted options={fontsize=\scriptsize}, listing only, colback=codebg, colframe=codebg, rounded corners}
// Функцию написали задолго до
// появления DoubleList
void fill(List *l, int val, size_t num) {
  for(size_t i = 0; i < num; i++)
  l->push_back(val);
}

// Код работает без изменений в
// функции fill
List l(5);
DoubleList dl(6);

fill(&l, 10, 42);
fill(&dl, 100, 24);
\end{tcblisting}
\end{itemize}
\end{frame}

\begin{frame}[fragile]
\frametitle{Пример: логирование}
\begin{itemize}
\usemintedstyle{tango}
\item Задача: хотим реализовать расширяемую систему логирования
\pause
\item Должна позволять логировать сообщения одновременно в разные места (т.н. \textit{sinks} -- в \mintinline{cpp}|stdout|, в файл, и т.п.)
\pause
\item При добавлении новых sink'ов не хочется переписывать весь код для их поддержки
\end{itemize}
\end{frame}

\begin{frame}[fragile]
\frametitle{Пример: логирование}
\begin{itemize}
\usemintedstyle{tango}
\item Хотим, чтобы код выглядел как-то так:
\usemintedstyle{lightbulb}
\begin{tcblisting}{listing engine=minted, minted language=cpp, minted options={fontsize=\scriptsize}, listing only, colback=codebg, colframe=codebg, rounded corners}
Logger logger;

StdoutSink stdoutSink(WARNING);
FileSink fileSink1("log.txt", WARNING);
FileSink fileSink2("debug_log.txt", DEBUG);
logger.addSink(stdoutSink);
logger.addSink(fileSink1);
logger.addSink(fileSink2);

logger.log(INFO, "Initialized");
logger.log(DEBUG, "Start time: ...");
\end{tcblisting}
\pause
\usemintedstyle{tango}
\item Хотим сделать так, чтобы при добавлении новых sink'ов класс \mintinline{cpp}|Logger| не пришлось переделывать
\end{itemize}
\end{frame}

\begin{frame}[fragile]
\frametitle{Пример: логирование}
\begin{itemize}
\usemintedstyle{tango}
\item Базовый класс:
\usemintedstyle{lightbulb}
\begin{tcblisting}{listing engine=minted, minted language=cpp, minted options={fontsize=\scriptsize}, listing only, colback=codebg, colframe=codebg, rounded corners}
class Sink {
protected:
  Level level_;
public:
  Sink(Level level): level_(level) {}
  virtual ~Sink() {}
  virtual void log(Level level, const char* message) {}
};
\end{tcblisting}
\end{itemize}
\end{frame}

\begin{frame}[fragile]
\frametitle{Пример: логирование}
\begin{itemize}
\usemintedstyle{tango}
\item Реализация логирования в \mintinline{cpp}|stdout|:
\usemintedstyle{lightbulb}
\begin{tcblisting}{listing engine=minted, minted language=cpp, minted options={fontsize=\scriptsize}, listing only, colback=codebg, colframe=codebg, rounded corners}
class StdoutSink : public Sink {
public:
  StdoutSink(Level level): Sink(level_) {}

  void log(Level level, const char* message) override {
    if (level >= level_)
      printf(message);
  }
};
\end{tcblisting}
\end{itemize}
\end{frame}

\begin{frame}[fragile]
\frametitle{Пример: логирование}
\begin{itemize}
\usemintedstyle{tango}
\item Реализация логирования в файл:
\usemintedstyle{lightbulb}
\begin{tcblisting}{listing engine=minted, minted language=cpp, minted options={fontsize=\scriptsize}, listing only, colback=codebg, colframe=codebg, rounded corners}
class FileSink : public Sink {
protected:
  FILE* file_;
public:
  FileSink(const char* filename, Level level)
    : Sink(level_)
    , file_(fopen(filename, "wt"))
  {}

  void log(Level level, const char* message) override {
    if (level >= level_)
      fprintf(file_, message);
  }

  ~FileSink() override {
    fclose(file_);
  }
};
\end{tcblisting}
\end{itemize}
\end{frame}

\begin{frame}[fragile]
\frametitle{Пример: логирование}
\begin{itemize}
\usemintedstyle{tango}
\item Реализация основного класса \mintinline{cpp}|Logger|:
\usemintedstyle{lightbulb}
\begin{tcblisting}{listing engine=minted, minted language=cpp, minted options={fontsize=\scriptsize}, listing only, colback=codebg, colframe=codebg, rounded corners}
class Logger {
private:
  Sink sinks_[100];
  size_t size_;
public:
  Logger(): size_(0) {}

  void addSink(Sink& sink) {
    sinks_[size_++] = sink;
  }

  void log(Level level, const char* message) {
    for (size_t i = 0; i < size_; ++i)
      sinks_[i].log(level, message);
  }
};
\end{tcblisting}
\pause
\usemintedstyle{tango}
\item Будет ли работать? \pause Нет, так как \mintinline{cpp}|Logger| хранит \textit{копии} переданных sink'ов, к тому же в виде объектов класса \mintinline{cpp}|Sink|
\end{itemize}
\end{frame}

\begin{frame}[fragile]
\frametitle{Пример: логирование}
\begin{itemize}
\usemintedstyle{tango}
\item При публичном наследновании объект класса-наследника неявно приводится к типу базового класса
\pause
\item \(\Longrightarrow\) Такой код скопилируется, хотя смысла в нём немного:
\usemintedstyle{lightbulb}
\begin{tcblisting}{listing engine=minted, minted language=cpp, minted options={fontsize=\scriptsize}, listing only, colback=codebg, colframe=codebg, rounded corners}
FileSink fileSink(...);
Sink sinkCopy = fileSink;
\end{tcblisting}
\pause
\usemintedstyle{tango}
\item Скопируется только часть объекта \mintinline{cpp}|fileSink|, отвечающая классу \mintinline{cpp}|Sink|
\item Вызов методов у \mintinline{cpp}|sinkCopy| вызовет методы класса \mintinline{cpp}|Sink|
\end{itemize}
\end{frame}

\begin{frame}[fragile]
\frametitle{Пример: логирование}
\begin{itemize}
\usemintedstyle{tango}
\item Исправление:
\usemintedstyle{lightbulb}
\begin{tcblisting}{listing engine=minted, minted language=cpp, minted options={fontsize=\scriptsize}, listing only, colback=codebg, colframe=codebg, rounded corners}
class Logger {
private:
  Sink* sinks_[100];
  size_t size_;
public:
  Logger(): size_(0) {}

  void addSink(Sink* sink) {
    sinks_[size_++] = sink;
  }

  ~Logger() {
    for (size_t i = 0; i < size_; ++i) {
      delete sinks_[i];
    }
  }
};
\end{tcblisting}
\end{itemize}
\end{frame}

\begin{frame}[fragile]
\frametitle{Pure virtual}
\begin{itemize}
\usemintedstyle{tango}
\item У реализации нового класса sink можно забыть реализовать метод \mintinline{cpp}|log()|, и тогда вызовется метод базового класса
\pause
\item Решение -- \textit{чисто виртуальный метод (pure virtual method)}:
\usemintedstyle{lightbulb}
\begin{tcblisting}{listing engine=minted, minted language=cpp, minted options={fontsize=\scriptsize}, listing only, colback=codebg, colframe=codebg, rounded corners}
class Sink {
public:
  ...
  virtual void log(Level level, const char* message) = 0;
};
\end{tcblisting}
\pause
\item Класс с чисто виртуальными методами называется \textit{абстрактным}
\pause
\item Создать экземпляр абстрактного класса \textbf{нельзя} (будет ошибка компиляции), но можно от него отнаследоваться
\pause
\item Если класс-наследник не реализует часть чисто виртуальных методов базового класса, то он останется абстрактным
\end{itemize}
\end{frame}

\begin{frame}[fragile]
\frametitle{Static vs dynamic dispatch}
\begin{itemize}
\usemintedstyle{tango}
\item Как работают виртуальные функции внутри, на уровне бинарного кода?
\pause
\item Если у объекта вызван невиртуальный метод, он вызывается как обычная функция \textit{(почти, см. thiscall)} -- это т.н. \textit{static dispatch}
\pause
\item Если функция виртуальная, нужно понять, какого типа на самом деле объект, и вызвать соответствующий метод -- это т.н. \textit{dynamic dispatch}
\pause
\item В C++ (и многих других ООП-языках) dynamic dispatch работает с помощью \textit{таблицы виртуальных функций}
\end{itemize}
\end{frame}

\begin{frame}[fragile]
\frametitle{Таблица виртуальных функций}
\begin{itemize}
\usemintedstyle{tango}
\item Для каждого класса с виртуальными методами компилятор создаёт особый массив указателей на функции
\pause
\item Для каждого виртуального метода где-то в этом массиве лежит указатель на него
\pause
\item В каждый экземпляр класса компилятор добавляет указатель на этот массив в некое фиксированное место (часто -- в начало класса)
\pause
\item При вызове виртуального метода сначала достаётся нужный элемент из этого массива, и потом вызывается как функция
\pause
\item \(\Longrightarrow\) Вызов виртуального метода немного дороже, чем вызов обычного метода
\end{itemize}
\end{frame}

\end{document}